\section*{Abstract}

Property-based testing tools such as Hypothesis can turn a formal grammar
into a generator of syntactically valid inputs, but the generator itself is
typically static: hand-tuned once and reused unchanged for the rest of a
fuzzing campaign. We study whether a large language model (LLM), given only
a formal grammar and parser-level feedback -- with no coverage
instrumentation -- can iteratively refine such a generator. We build a fully
black-box pipeline targeting the cJSON parser (pinned at v1.7.19): a baseline
Hypothesis strategy derived directly from an ANTLR JSON grammar generates
inputs, a sanitizer-instrumented harness classifies each as accepted,
rejected, or crashing, and a refinement loop feeds an LLM a summary of three
observable proxy signals -- acceptance rate, the structural diversity of the
shapes produced, and the distinct rejection signatures seen -- asking it to
propose a revised generator. Every proposal is statically sandboxed
(AST-checked, import-restricted) and validated before it ever touches the
target binary. Across five independent, seeded runs of five refinement
iterations each, LLM-guided refinement raises mean acceptance rate from
57.6\% to 96.8\% and structural diversity by 123\% relative to a static
baseline, using no coverage data at all. We additionally run the identical,
unmodified pipeline against a second parser (parson) for five real
iterations (3{,}000 sanitizer-instrumented executions); no crash was found,
and we report this as a negative result, supported by manual review of the
highest-risk code paths, rather than claim a bug that was not observed. We
discuss how grammar-implementation mismatches -- for instance, duplicate
JSON object keys accepted by one parser and rejected by the other under an
identical grammar -- are themselves a source of fuzzing-relevant signal that
a purely formal grammar cannot supply, and we release the full pipeline,
harnesses, and reproducibility tooling.
